\begin{center}
\textit{Genera actually correspond to all characters except those found impossible in Articles $263,264$.}
\end{center}

\Needspace{8\baselineskip}
\Article{287}
We add the following further remarks on the problem solved in the preceding article.

I.\ If the form $F$ is transformed by the substitution $p,p',p'',p''';\ q,q',q'',q'''$ into
the product of two forms $(h,i,k)$, $(h',i',k')$ (both, as we always
\SourceLeaf{342}
suppose, understood in the proper sense), we have the following equations,
which are readily derived from conclusion 3 of Article 235:
\[
\begin{gathered}
p''hn'-p'h'n-p(in'-i'n)=0\\
(p''-p')(in'+i'n)-p(kn'-k'n)+p'''(hn'-h'n)=0\\
p'kn'-p''k'n-p'''(in'-i'n)=0
\end{gathered}
\]
and three others obtained by interchanging the numbers $p$, $p'$,
$p''$, $p'''$ with $q$, $q'$, $q''$, $q'''$. Here $n$,
$n'$ are the positive square roots of the quotients obtained by
dividing the determinants of the forms $(h,i,k)$, $(h',i',k')$ by the determinant
of the form $F$. So, if these forms are identical, or $n=n'$,
$h=h'$, $i=i'$, $k=k'$, those equations become
\[
(p''-p')hn=0,\quad (p''-p')in=0,\quad (p''-p')kn=0
\]
This \textit{necessarily} gives $p'=p''$, and in exactly the same way
$q'=q''$. --- So, giving the forms $(h,i,k)$, $(h',i',k')$ the \textit{same}
indeterminates $t$, $u$, and denoting the indeterminates of the
form $F$ by $T$, $U$, the substitution transforms $F$ as follows:
\[
T=ptt+2p'tu+p'''uu,\quad U=qtt+2q'tu+q'''uu\quad\text{into}\quad (htt+2itu+kuu)^2
\]

II.\ If the form $F$ arises by doubling the form $f$, it also
arises by doubling any other form in the same class as $f$; in other words, the
class of the form $F$ arises by doubling the class of the form $f$
(see Article 238). Thus, in the example of the preceding article, $(5,2,31)$
also arises by doubling the form $(11,-5,16)$, which is properly equivalent
to $(11,17,40)$. From one class whose doubling gives the class of the form
$F$, we find \textit{all} such classes (if there are more) with the
help of Problem 260. In our example there is no other positive class of
this kind, since there is only one positive properly primitive ambiguous
class of determinant $-151$ (namely the principal class). Composing the
only negative ambiguous class $(-1,0,-151)$ with the class $(11,-5,16)$ gives the
class $(-11,-5.\ -16)$, so this will be the only negative class whose doubling
gives the class $(5,2,31)$.

III.\ The very solution of the problem in the preceding article has
shown that every properly primitive (positive) class of binary forms
belonging to the principal genus can arise by doubling some properly
primitive class of the same determinant. So the theorem of Article 261,
which assured us that \textit{at least} half of all characters
\SourceLeaf{343}
assignable to a given nonsquare determinant $D$ cannot correspond
to properly primitive (positive) genera, can now be extended: such genera
actually correspond to \textit{exactly} half of all these characters,
and so none can correspond to the other half (see the proof of that
theorem). In Article 264 all these assignable characters were divided
equally into two kinds $P$, $Q$. The latter $Q$ were proved
unable to correspond to properly primitive (positive) forms, while it
remained uncertain whether genera always actually correspond to each
of the remaining $P$. This doubt is now completely removed, and
we are sure that in the whole collection of characters $P$ there
is none without a corresponding genus. --- It also follows readily
that, for a negative determinant, in the \textit{negative} properly
primitive order, where Article 264, I showed all $P$ to be impossible
and \textit{only} $Q$ possible, \textit{all} $Q$ are actually
possible. For let $K$ denote any character in $Q$, $f$ any
form in the negative properly primitive order of forms of determinant
$D$, and $K'$ its character, which will be in $Q$. Then the
character composed from $K$, $K'$ (according to Article 246)
clearly belongs to $P$. So there are positive properly primitive
forms of determinant $D$ corresponding to it. Composing such a form
with $f$ clearly gives a negative properly primitive form of determinant
$D$ whose character is $K$. --- Exactly the same argument proves
that, in the improperly primitive order, the characters found by the
rules of Article 264 II, III to be the \textit{only} possible ones are
\textit{all} possible, whether they are $P$ or $Q$. --- Unless
we are very much mistaken, these theorems rank among the most beautiful
in the theory of binary forms, all the more so because, although they
are extremely simple, they lie so deep that a rigorous proof cannot be
given without the help of so many other investigations.

\Needspace{8\baselineskip}
\begin{center}
\textit{Theory of decomposing both numbers and binary forms into three squares.}
\end{center}
We now turn to another application of the preceding digression: the
decomposition of both numbers and binary forms into three squares.
We begin with the following

\Needspace{8\baselineskip}
\Article{288}
\textsc{Problem.} \textit{Let $M$ denote a positive number. Find the
conditions under which there can be primitive negative binary forms of
determinant $-M$ that are quadratic residues of $M$, that is, for which
$1$ is a characteristic number.}

\SourceLeaf{344}
\textit{Solution.} Let $\Omega$ denote the collection of all particular
characters given by the relations of the number $1$ both to each
(odd) prime divisor of $M$ and to the number $8$ or $4$
when it divides $M$. These characters will clearly be $Rp$,
$Rp'$, $Rp''$, and so on, where $p$, $p'$, $p''$, and so on
denote those prime divisors; also $1,4$ when $4$ divides it, and
$1,8$ when $8$ divides $M$. Further, let us use the letters
$P$, $Q$ in the same sense as in the preceding article, or in
Article 264. We now distinguish the following cases.

I.\ When $M$ is divisible by $4$, $\Omega$ will be a complete
character. Article 233 V shows that $1$ can be a characteristic
number only of forms whose character is $\Omega$. But $\Omega$ is
clearly the character of the principal form $(1,0,M)$, so it belongs to
$P$ and therefore cannot belong to a properly primitive negative
form. Since there are no improperly primitive forms for such a determinant,
in this case there are no primitive negative forms at all that are
residues of $M$.

II.\ When $M\equiv3\ (\mathrm{mod.}\ 4)$, exactly the same arguments hold, with the sole
exception that in this case there is an \textit{improperly} primitive
negative order. In it the characters $P$ will be possible or
impossible according as $M\equiv3$ or $\equiv7\ (\mathrm{mod.}\ 8)$; see Article 264, III.
So, in the first case there will be a genus in this order whose character
is $\Omega$, and $1$ will be a characteristic number of all forms
in it. In the second case there can be no negative forms at all with
this property.

III.\ When $M\equiv1\ (\mathrm{mod.}\ 4)$, $\Omega$ is not yet a complete character: the
relation to the number $4$ must be added. But $\Omega$ clearly
must enter the character of a form whose characteristic number is
$1$; conversely, any form whose character is either $\Omega;\ 1,4$ or
$\Omega;\ 3,4$ has characteristic number $1$. Now $\Omega;\ 1,4$ is clearly
the character of the principal genus, which belongs to $P$, and
so is impossible in the negative properly primitive order. For the same
reason, $\Omega;\ 3,4$ will belong to $Q$ (Article 263), so a genus in
the negative properly primitive order will correspond to it, all of
whose forms will have characteristic number $1$. There is no
improperly primitive order in this case, or in the next.

IV.\ When $M\equiv2\ (\mathrm{mod.}\ 4)$, we must add to $\Omega$ the relation to $8$
to make a complete character: either $1\mathbin{\textit{and}}3,8$ or $5\mathbin{\textit{and}}7,8$ when $M\equiv2\ (\mathrm{mod.}\ 8)$,
and either $1\mathbin{\textit{and}}7,8$ or $3\mathbin{\textit{and}}5,8$ when $M\equiv6\ (\mathrm{mod.}\ 8)$. In the first case the
character $\Omega;\ 1\mathbin{\textit{and}}3,\allowbreak 8$ clearly belongs to $P$, and so $\Omega;\ 5\mathbin{\textit{and}}7,\allowbreak 8$ to
$Q$. Thus a negative properly primitive genus will
\SourceLeaf{345}
correspond to it. In the second case, the same argument gives one genus
in the negative properly primitive order whose forms have the prescribed
property, namely the one whose character is $\Omega;\ 3\mathbin{\textit{and}}5,\allowbreak 8$.

It follows that primitive negative forms of determinant $-M$ with
characteristic number $1$ exist when $M$ is congruent to one
of the numbers $1$, $2$, $3$, $5$, $6$ modulo
$8$, and always in just one genus, which will be improper when
$M\equiv3$. There are no such forms at all when $M\equiv0,\allowbreak \ 4\text{ or }7\ (\mathrm{mod.}\ 8)$. Further, if
$(-a,-b,-c)$ is a primitive negative form with characteristic number
$+1$, clearly $(a,b,c)$ is a primitive positive form with characteristic
number $-1$. So, in the first five cases (when $M\equiv1,2,3,5,6$), there is
one primitive positive genus whose forms have characteristic number
$-1$, and it is \textit{improper} for $M\equiv3$. In the remaining
three cases (when $M\equiv0,4,7$) there can be no such positive forms at all.

